\subsection{An explicit inertia witness at seven}
The same strict reduction also has a local inertia construction.
\begin{proposition}\label{prop:seven-inertia-mixed-prime-witness}
There is an arithmetic element \(w\in T_D\) for which every
single-prime word action is trivial and
\[
 W_7(w)=\left(\frac{1+7i}{1-7i}\right)^{-2}.
\]
This value has infinite order.
\end{proposition}
\begin{proof}
Fix decomposition groups at the primes two and seven,
with inertia groups \(I_2,I_7\subset H\).  The normalized
tangential action on the free pro-\(p\) fundamental group of
\(\mathbf P^1-\{0,1,\infty\}\) is unramified outside \(p\).
This is a statement about the tangential action, not just the
outer action; see \cite[\S8.5]{HainMatsumoto2003}.  Thus
\begin{equation}\label{eq:disjoint-inertia-word-actions}
                  \Phi_p(I_r)=1\qquad(r\ne p).
\end{equation}

Here is the local CM input in an explicit form.  Let
\(K=\mathbf Q_7(i)\) and \(O=\mathbf Z_7[i]\).
This is the unramified quadratic extension of \(\mathbf Q_7\).
The curve \(E_4\) has good reduction at seven, and a direct
count gives \(\#E_4(\mathbf F_7)=8\): the numbers of affine
points above \(x=0,\ldots,6\) are respectively
\(1,1,0,0,2,2,1\), with the point at infinity added.
Consequently the seven-power Frobenius \(\pi\) on the
reduction satisfies \(\pi^2=[-7]\).
In the formal parameter \(t=-x/y\), the endomorphism
\([-7]\) of \(\widehat E_4/O\) therefore satisfies
\[
 [-7](t)\equiv-7t\pmod{t^2},\qquad
 [-7](t)\equiv t^{49}\pmod7.
\]
The CM automorphism \((x,y)\mapsto(-x,iy)\) acts on this
parameter by \(t\mapsto it\).  Its action and the
\(\mathbf Z_7\)-action give the formal \(O\)-module structure.
The preceding two congruences identify this formal module
with the Lubin--Tate module for the uniformizer \(-7\).
Since \([7]\) is purely inseparable on the reduction, every
point of \(E_4[7^n]\) reduces to the identity.  Thus its
Tate module is the rank-one \(O\)-module of this formal
group.  The Lubin--Tate reciprocity theorem says that the
inertia image on it is all of \(O^\times\); see
\cite[p.~380, the displayed reciprocity formula~(2)]{LubinTate1965}.
Indeed the finite torsion fields are totally ramified with
groups \((O/7^nO)^\times\), and inverse limits give the
assertion.  As \(K/\mathbf Q_7\) is unramified, its inertia
is \(I_7\).

Choose
\[
             a=\frac{1+7i}{1-7i}\in1+7O.
\]
Then \(a\bar a=1\), and \(a\) has infinite order.  For
example its seven-adic logarithm has first term \(14i\),
of valuation one, while all subsequent terms have larger
valuation.  Choose \(j\in I_7\) whose action on \(T_7(E_4)\)
is multiplication by \(a\).  The Weil pairing gives
\(\lambda_7(j)=\det(a)=a\bar a=1\).  Inertia at seven
fixes all roots of unity of order prime to seven, so in fact
\begin{equation}\label{eq:inertia-seven-full-multiplier}
                              \lambda(j)=1.
\end{equation}

The local cyclotomic extension at two gives an element
\(h\in I_2\) with \(\lambda_2(h)=-1\).
All other components of \(\lambda(h)\) are one.
Since \(h(i)=-i\), its action on the seven-adic CM module
is semilinear for the involution of \(O\).  Whatever its
unit factor, it therefore conjugates multiplication by
\(a\) to multiplication by \(\bar a\).

Use the convention \([h,j]=h j h^{-1}j^{-1}\), and put
\(z=[h,j]\).  Its full multiplier is one.  For each prime
\(p\), at least one of \(\Phi_p(h),\Phi_p(j)\) is the
identity by \eqref{eq:disjoint-inertia-word-actions}, so
\begin{equation}\label{eq:commutator-all-word-kernels}
                           \Phi_p(z)=1\quad\hbox{for every }p.
\end{equation}
For any rational-prime Kummer cocycle \(\rho_q\), the
crossed-homomorphism law and \(\lambda(j)=1\) give
\[
              \rho_q(z)=(\lambda(h)-1)\rho_q(j).
\]
This adelic vector is supported only at its two-component.
Let \(e_7\in\widehat{\mathbf Z}=\prod_p\mathbf Z_p\)
have seven-component one and all other components zero,
and define \(w=z^{e_7}\).  The exponent is taken through
the continuous map
\(\widehat{\mathbf Z}\to\overline{\langle z\rangle}\);
it is a profinite group power.  On the multiplier-one
closed cyclic group generated by \(z\), all Kummer
cocycles are additive.  Hence
\[
                  \rho_q(w)=e_7\rho_q(z)=0\quad\hbox{for all }q.
\]
Thus \(w\in D\cap H\), and
\eqref{eq:commutator-all-word-kernels} implies
\(w\in P_D\cap H\).  The cubic residue defect is zero
on arithmetic elements, so \(w\in T_D\cap H\).

Finally, in the signed CM identification,
\[
                    W_7(z)=\bar a/a=a^{-2}.
\]
This element belongs to the pro-seven group \(1+7O\),
so exponentiation by \(e_7\) leaves it unchanged:
\[
                         W_7(w)=a^{-2}.
\]
Projection from the local norm-one torus to its pro-seven
Sylow factor annihilates every pro-two element.  It
therefore factors through \(N_4(4)_{\rm odd}\).
The displayed value has infinite order, so the image of
the closed cyclic subgroup generated by \(w\) in
\(T_D/U_D\) is infinite.  This proves both strictness
assertions.
\end{proof}
